\documentclass[french,a4paper]{article}

\usepackage[utf8]{inputenc}
\usepackage[T1]{fontenc}
\usepackage{lmodern}
\usepackage{geometry}
\usepackage{graphicx}
\usepackage{babel}
\usepackage{amsmath}
\usepackage{amsthm}
\usepackage{amssymb}
\usepackage{hyperref}
\usepackage{stmaryrd}
\usepackage{enumitem}
\usepackage{tcolorbox}
\usepackage{dsfont}
\usepackage{multirow}
\usepackage{etoc}
\usepackage{titlesec}
\usepackage{esint}
\usepackage{cancel}
\usepackage{mathdots}

\setlist[itemize]{label=$\bullet$}


\renewcommand{\P}{\mathbb{P}}
\newcommand{\N}{\mathbb{N}}
\newcommand{\Z}{\mathbb{Z}}
\newcommand{\Q}{\mathbb{Q}}
\newcommand{\R}{\mathbb{R}}
\newcommand{\C}{\mathbb{C}}
\newcommand{\K}{\mathbb{K}}
\renewcommand{\L}{\mathbb{L}}
\newcommand{\U}{\mathbb{U}}
\newcommand{\st}{^*}
\newcommand{\tq}{\middle|}
\newcommand{\inv}{^{-1}}
\newcommand{\E}{\mathbb{E}}
\newcommand{\F}{\mathbb{F}}
\newcommand{\V}{\mathbb{V}}
\renewcommand{\ker}{\text{Ker}}
\newcommand{\im}{\text{Im}}
\newcommand{\card}{\text{Card}}
\newcommand{\Tr}{\text{Tr}}

\theoremstyle{plain}
\newtheorem{thm}{Théorème}
\newtheorem*{ind}{Indication}

\theoremstyle{definition}
\newtheorem{dfn}{Definition}

\theoremstyle{remark}
\newtheorem*{rmq}{Remarque}
\newtheorem*{app}{Application}

\newtcolorbox[auto counter]{ex}[2][]{colback=white, colframe=green!50!white, coltitle=black, fonttitle=\bfseries, title={Exercice~\thetcbcounter~: #2}, after title={\hfill #1}}

\title{TD Fonctions numériques et vectorielles d'une variable réelle}

\begin{document}

\maketitle

Dans tous les exercices, $E$ désigne un espace vectoriel normé de dimension finie. Les fonctions considérées sont définies sur un intervalle de $\R$ et à valeurs dans $E$. pour certains exercices, on se limitera à $E=\K$, avec $\K=\R$ ou $\K=\C$ (fonctions numériques).
\section{Dérivation des fonctions à valeurs réelles}

\begin{ex}[Moulin.FV.01]{$\Delta$ Théorème de Darboux}
	Montrer que la dérivée d'une fonction réelle dérivable sur un intervalle possède la propriété des valeurs intermédiaires :
	\begin{itemize}
		\item elle ne peut changer de signe sans s'annuler ;
		\item l'image directe de n'importe quel sous-intervalle de l'intervalle de départ est un intervalle.
	\end{itemize}
\end{ex}
\begin{proof}
	Supposer sans perte de généralité qu'il existe $a<b$ tels que $f'(a)>0$ et $f'(b)<0$ puis se ramener à l'utilisation du théorème de Rolle. Pour la deuxième propriété, utiliser la première et $x\mapsto f(x)-tx$ si on prend $f'(a)<t<f'(b)$.
\end{proof}

\begin{ex}[Moulin.FV.02]{L'image d'un segment par la dérivée est $\R$}
	Trouver une fonction dérivable $f:[0,1]\to\R$ telle que $f'([0,1])=\R$.
\end{ex}
\begin{proof}
	Considérer par exemple $$x\mapsto x^{3/2}\sin\left(\frac{1}{x}\right)$$ prolongée par continuité par $0$ en $0$. On montre qu'elle est bien dérivable et qu'elle vérifie la propriété demandée (elle oscille de plus en plus vite).
\end{proof}

\begin{ex}[Moulin.FV.03]{$\circ$}
	Soit $f:[a,b]\to\R$ dérivable telle que : $$f'(a)>0\ \ f'(b)>0\ \ \text{et}\ \ f(a)=f(b)=0$$ Montrer qu'il existe $c\in]a,b[$ tel que $f(c)=0$ et $f'(c)\leqslant 0$.
\end{ex}

\begin{ex}[Moulin.FV.04]{$\Delta$ Classique}
	Soit $f:[a,b]\to\R$ dérivable. Montrer qu'il existe $c\in]a,b[$ tel que $\displaystyle f'(x)=\frac{f(c)-f(a)}{c-a}$ dans chacun des deux cas suivants : 
	\begin{itemize}
		\item $f(a)=f(b)=0$ et $f'(a)=0$
		\item $f'(b)=f'(a)$
	\end{itemize}
\end{ex}
\begin{proof}
	On considère $$g(x)=\int_{0}^{1}f'\left(\frac{t-a}{b-a}x+\frac{b-t}{b-a}a\right)\ dt$$ qui vaut $\displaystyle\frac{f(x)-f(a)}{x-a}$ si $x\in]a,b]$ et $f'(a)$ si $x=a$ (si je ne me suis pas trompé dans le calcul de tête, ce qui est fort peu probable).
	\begin{itemize}
		\item Lui appliquer le théorème de Rolle pour obtenir le résultat.
		\item Si $g$ admet un extremum en $c\in]a,b[$, conclure comme précédemment. Sinon, on suppose sans perte de généralité : $$\forall x\in[a,b],\ g(a)\leqslant g(x)\leqslant g(b)$$ On en déduit que $$\forall x\in]a,b],\ f(x)\leqslant f(a)+\frac{f(b)-f(a)}{b-a}(x-a)$$ En calculant, on en déduit que $$\forall x\in]a,b],\ \frac{f(b)-f(x)}{b-x}\geqslant g(b)$$ En passant à la limite lorsque $x$ tend vers $b$, on obtient $f'(b)\geqslant g(b)$. On en déduit que $g(a)=g(b)$ et on conclut comme précédemment.
	\end{itemize}
\end{proof}

\begin{ex}[Moulin.FV.05]{}
	Soit $\lambda_1,\ldots,\lambda_n$ des réels deux à deux distincts, $P_1,\ldots,P_n$ des polynômes réels non nuls et $f:\R\to\R$ définie par : $$\forall x\in\R,\ f(x)=\sum_{k=1}^{n}e^{\lambda_k x}P_k(x)$$ Montrer que $f$ s'annule au plus $n-1+d$ fois, où $d=\displaystyle\sum_{k=1}^{n}\deg(P_k)$.
\end{ex}

\begin{ex}[Moulin.FV.06]{}
	Soit $f$ une application de $\R$ dans $\R_+$ de classe $\mathcal{C}^1$ admettant une limite nulle en $+\infty$. Montrer qu'il existe une suite $(x_n)_{n\in\N}$ de réels telle que la suite $(f'(x_n))_{n\in\N}$ converge vers $0$.
\end{ex}
\begin{proof}
	Construire par récurrence une telle suite en utilisant le théorème des accroissements finis.
\end{proof}

\begin{ex}[?]{Une étrange équation}
	Résoudre l'équation $5^x+2^x=4^x+3^x$ d'inconnue $x\in\R$.
\end{ex}
\begin{proof}
	Considérer $f:t\mapsto t^x$ sur $\R_+\st$ qui est dérivable de dérive $f':t\mapsto xt^{x-1}$. L'équation se réécrit $f(5)-f(4)=f(3)-f(2)$. Par le théorème des accroissements finis, cela se réécrit $$f'(t_1)(5-4)=f'(t_2)(3-2)$$ pour un certain $2<t_2<3$ et un certain $4<t_1<5$. Cela se réécrit donc $xt_1^{x-1}=xt_2^{x-1}$ donc nécessairement $x=1$ car $t_1\ne t_2$. Et réciproquement, $x=1$ est bien solution.
\end{proof}

\section{Dérivation}

\begin{ex}[Moulin.FV.07]{$\circ$}
	Soit $A$ une application dérivable de $\R$ dans $\mathcal{M}_n(\R)$. Prouver l'équivalence entre les propositions suivantes :
	\begin{itemize}
		\item $\forall (s,t)\in\R^2,\ A(s)A(t)=A(t)A(s)$ 
		\item $\forall (s,t)\in\R^2,\ A(s)A'(t)=A'(t)A(s)$
	\end{itemize}
\end{ex}
\begin{proof}
	Dans le sens direct, dériver par rapport à $t$. Dans le sens réciproque, intégrer.
\end{proof}

\begin{ex}[Moulin.FV.08]{$\Delta$ Dérivée des puissances entières}
	Soit $A$ une application dérivable d'un intervalle $I$ de $\R$ dans $\mathcal{M}_n(\R)$.
	\begin{enumerate}
		\item Quelle est la dérivée de $A^p$ pour $p\in\N$ ?
		\item On suppose $A(t)$ inversible pour tout $t\in I$. Quelle est la dérivée de $A^{-p}$ pour $p\in\N\st$ ?
	\end{enumerate}
\end{ex}
\begin{proof}
	Attention aux idées reçues !
	\begin{enumerate}
		\item \textit{A priori}, ce n'est pas $pA'A^{p-1}$ car on ne sait pas si $A$ commute avec $A'$. Le théorème de dérivation d'une application multilinéaire assure qu'il s'agit plutôt de : $$\sum_{k=1}^pA^{k-1}A'A^{p-k}$$
		\item Écrire $AA\inv=I_n$ et dériver cette relation. On obtient : $$(A\inv)'=-A\inv A' A\inv$$ Pour les puissances négatives, utiliser la question $1$ et combiner avec le résultat précédent.
	\end{enumerate}
\end{proof}

\begin{ex}[Moulin.FV.09]{$\Delta$ Utilisation d'une série de fonctions}
	Soit $f:[0,1]\to E$ continue en $0$. On suppose que $t\mapsto\displaystyle\frac{f(3t)-f(t)}{t}$ admet en $0$ une limite $l\in E$. Montrer que $f$ est dérivable en $0$ et exprimer $f'(0)$ en fonction de $l$. On pourra utiliser une série de fonctions.
\end{ex}
\begin{proof}
	Écrire $$f(t)-f(0)=\sum_{n=0}^{+\infty}\left[f\left(\frac{t}{3^{n}}\right)-f\left(\frac{t}{3^{n+1}}\right)\right]$$ On en déduit que : $$\frac{f(t)-f(0)}{t}=\sum_{n=0}^{+\infty}\frac{1}{3^{n+1}}\left[\frac{\displaystyle f\left(3\times\frac{t}{3^{n+1}}\right)-f\left(\frac{t}{3^{n+1}}\right)}{\displaystyle\frac{t}{3^{n+1}}}\right]$$ Or, on a une convergence normale donc par interversion de limites $f$ est dérivable et $$f'(0)=\sum_{n=0}^{+\infty}\frac{l}{3^{n+1}}=\frac{l}{2}$$
\end{proof}

\begin{ex}[Moulin.FV.10]{$\circ$}
	Soit $I$ un intervalle de $\R$ non trivial et $O:I\to\mathcal{O}_n(\R)$. Montrer qu'il existe une fonction $A:I\to\mathcal{A}_n(\R)$ telle que $$\forall t\in I,\ O'(t)=O(t)A(t)$$
\end{ex}
\begin{proof}
	On considère $O^T$ qui est elle aussi dérivable par linéarité de la transposition et à valeurs dans $\mathcal{O}_n(\R)$. Elle vérifie $(O^T)'=(O')^T$ par linéarité. Alors : $$\forall t\in I,\ O(t)O^T(t)=I_n$$ En dérivant cette relation, on obtient : $$\forall t\in I,\ O'(t)O^T(t)+O(t)O'^T(t)=0$$ Ainsi : $$\forall t\in I,\ O'(t)=O(t)\left[\underbrace{-(O\inv)'(t)O(t)}_{:=A(t)}\right]$$ Or, d'après un exercice précédent, on a : $$A(t)=O\inv(t)O'(t)O\inv(t)O(t)=O\inv(t)O'(t)=O^T(t)O'(t)$$ Or, on sait aussi que $$A(t)=-(O')^T(t)O(t)$$ Par conséquent : $$\forall t\in T,\ A(t)^T=-A(t)$$ donc $A$ est bien à valeurs antisymétriques.
\end{proof}

\begin{ex}[Moulin.FV.11]{Hélices}
	Soit $E$ un espace euclidien et $v\in E$ non nul. Déterminer les fonction dérivables $f:\R\to E$ telles que $(f'\mid v)$ soit une fonction constante.
\end{ex}
\begin{proof}
	Cela revient à avoir $(f\mid v)'=C$ avec $C$ une constante réelle. Intégrer puis voir ce que cela implique pour $f$. Penser à la physique et à la force de Lorentz : on obtient une hélice généralisée.
\end{proof}

\begin{ex}[Moulin.FV.12]{Dérivation du polynôme caractéristique}
	Si $M\in\mathcal{M}_n(\C)$, on note $\chi_M$ l'application : $$\begin{array}{rcl}
		\R&\longrightarrow&\C\\
		t&\longmapsto&\det(tI_n-M)
	\end{array}$$ Étant donnée $A\in\mathcal{M}_n(\C)$ et $i\in\llbracket1,n\rrbracket$, on note $A_i$ la matrice extraite de $A$ obtenue en supprimant la ligne et la colonne d'indices $i$.
	\begin{enumerate}
		\item Exprimer $\chi_A'$ en fonction des $\chi_{A_i}$.
		\item En déduire le coefficient en $X$ dans le polynôme caractéristique de $A$.
	\end{enumerate}
\end{ex}
\begin{proof}
	\begin{enumerate}
		\item Utiliser la règle de dérivation d'une forme multilinéaire et reconnaître l'apparition des déterminants des matrices $A_i$. On trouve : $$\chi_A'=\sum_{k=1}^{n}\chi_{A_k}$$
		\item Il correspond au coefficient constant de $\chi_A'$. C'est donc : $$(-1)^{n-1}\sum_{k=1}^{n}\det(A_k)$$
	\end{enumerate}
\end{proof}

\begin{ex}[Moulin.FV.13]{Dérivabilité de l'inverse dans une algèbre}
	Soit $A$ une algèbre de dimension finie et $f\in\mathcal{C}^1(I,A)$ à valeurs dans l'ensemble $\Omega$ des éléments inversibles de $A$. On définit $g$ sur $I$ par $g(t)=f\inv(t)$.
	\begin{enumerate}
		\item On suppose $g$ de classe $\mathcal{C}^1$. Exprimer sa dérivée en fonction de $f'$ (attention, on n'a pas supposé $A$ commutative).
		\item On suppose $A=\mathcal{M}_n(\K)$. Montrer que $g$ est de classe $\mathcal{C}^1$.
		\item $\star$ Démontrer que $g$ est de classe $\mathcal{C}^1$ dans le cas général. on pourra considérer l'application : $$\begin{array}{lrcl}
			\varphi:&A&\longrightarrow&\mathcal{L}(A) \\
			&a&\longmapsto&(x\mapsto ax)
		\end{array}$$
	\end{enumerate}
\end{ex}
\begin{proof}
	La dernière question est difficile.
	\begin{enumerate}
		\item Écrire $fg=1$ puis dériver cette relation pour obtenir $g'=-gf'g$.
		\item Utiliser la formule de la comatrice (tout ce qui la compose est $\mathcal{C}^1$).
		\item Il semblerait que $\varphi$ puisse nous permettre de nous ramener à la question 2.
	\end{enumerate}
\end{proof}

\section{Fonctions de classe $\mathcal{C}^k$}

\begin{ex}[Moulin.FV.14]{$\Delta$ Fonction de décollage en douceur}
	Montrer que la fonction : $$\begin{array}{lrcl}
		\varphi:&\R&\longrightarrow&\R \\
		&x&\longmapsto&\displaystyle\begin{cases}
			e^{-1/x}&\text{si }x>0 \\
			0&\text{sinon}
		\end{cases}
	\end{array}$$ est de classe $\mathcal{C}^\infty$.
\end{ex}
\begin{proof}
	Le seul problème est en $0$ à droite. Prouver par récurrence l'existence d'un polynôme qui apparaît devant l'exponentielle pour les dérivées.
\end{proof}
\begin{rmq}[Un contre-exemple classique]
	On a ici l'exemple d'une fonction de classe $\mathcal{C}^\infty$ sur $\R$ qui n'est pas pour autant développable en série entière : en effet, si cette fonction était développable en série entière, comme elle est nulle sur $]-1,0]$, elle serait nulle partout.
\end{rmq}

\begin{ex}[Moulin.FV.15]{}
	Soit $a<b$.
	\begin{enumerate}
		\item Trouver une fonction $f$ positive de classe $\mathcal{C}^\infty$ sur $\R$ telle que : $$\forall x\in\R,\ f(x)=0\iff x\notin]a,b[$$
		\item Trouver une fonction de classe $\mathcal{C}^\infty$ sur $\R$ nulle sur $]-\infty,a]$ et égale à $1$ sur $[b,\infty[$.
	\end{enumerate}
\end{ex}
\begin{proof}
	Utiliser les fonctions de décollage en douceur (exercice précédent) en les recollant convenablement.
\end{proof}

\begin{ex}[Moulin.FV.16]{$\circ$}
	Pour $\lambda\in\K$ et $k\in\N$, décrire $\ker(D-\lambda\text{id})^k$, où $D$ désigne la dérivation sur $\mathcal{C}^\infty(I,\K)$.
\end{ex}

\begin{ex}[Moulin.FV.17]{$\Delta$ Racines des dérivées successives de $\arctan$}
	On considère $f$ définie par $\forall t\in\R,\ f(t)=\displaystyle\frac{1}{1+t^2}$. 
	\begin{enumerate}
		\item Montrer que pour tout $n\in\N$, la dérivée $n$-ème de $f$ admet exactement $n$ racines réelles.
		\item Expliciter ces racines.
	\end{enumerate} 
\end{ex}
\begin{proof}
	Exercice classique. A faire...
\end{proof}

\begin{ex}[Moulin.FV.18]{$\Delta$ Théorème du relèvement}
	\begin{enumerate}
		\item Soit $f:I\mapsto\C\st$ de classe $\mathcal{C}^1$. Montrer qu'il existe $g:I\mapsto\C$ de classe $\mathcal{C}^1$ telle que $f=\exp\circ g$. Étudier dans quelle mesure $g$ est déterminée de manière unique par $f$.
		\item Montrer que si $f$ est de classe $\mathcal{C}^n$ avec $n\in\N\st\cup\{+\infty\}$ alors une telle fonction $g$ est de classe $\mathcal{C}^n$.
		\item Application : déterminer toutes les applications $f:\R\to\C$ de classe $\mathcal{C}^1$ telles que $$\forall t\in\R,\ f(t)^2=e^{it}$$
	\end{enumerate}
\end{ex}
\begin{proof}
	Exercice classique. A faire...
	\begin{enumerate}
		\item Unicité à $2ik\pi$ près.
		\item 
		\item Utiliser la première question.
	\end{enumerate}
\end{proof}

\begin{ex}[Moulin.FV.19]{$\Delta$ Une inégalité d'interpolation}
	Soit $f:[a,b]\to\R$ dérivable à l'ordre $n+1$. on fixe $a_0<a_1<\ldots<a_n$ dans $[a,b]$ et on considère le polynôme $P\in\R_n[X]$ interpolant $f$ aux points $a_0,\ldots,a_n$.
	\begin{enumerate}
		\item $\star$ Montrer : $$\forall x\in[a,b],\ \exists t\in[a,b]\ :\ f(x)-P(x)=\frac{f^{(n+1)}(t)}{(n+1)!}\prod_{k=0}^{n}(x-a_k)$$
		\item En déduire : $$\mathcal{N}_\infty(f-P)\leqslant\frac{1}{(n+1)!}\mathcal{N}_\infty(f^{(n+1)})\underset{x\in[a,b]}{\sup}\prod_{k=0}^{n}\lvert x-a_k\rvert$$
	\end{enumerate}
\end{ex}
\begin{proof}
	La première question est difficile.
	\begin{enumerate}
		\item A faire...
		\item Conséquence directe de la première question avec des majorations brutales.
	\end{enumerate}
\end{proof}

\begin{ex}[Moulin.FV.20]{}
	Soit $f\in\mathcal{C}^{\infty}([0,+\infty[,\R)$ telle que $f(0)f'(0)>0$ et $\lim_{+\infty}f=0$. Montrer qu'il existe une suite strictement croissante $(x_n)_{n\in\N\st}$ telle que $\forall n\in\N\st,\ f^{(n)}(x_n)=0$.
\end{ex}

\section{Formules de Taylor}

\begin{ex}[Moulin.FV.21]{Formule des rectangles médians}
	Soit $f:[-1,1]\to\C$ une fonction de classe $\mathcal{C}^2$. Montrer que $$\left\lvert\int_{-1}^{1}f(t)\ dt-2f(0)\right\rvert\leqslant\frac{1}{3}\underset{x\in[-1,1]}{\sup}\lvert f''(x)\rvert$$
\end{ex}
\begin{proof}
	Appliquer la formule de Taylor avec reste intégral à l'ordre $3$ à $F$ une primitive de $f$ entre $-1$ et $0$ puis entre $1$ et $0$ et les sommer. Finir en majorant délicatement le polynôme dans le reste.
\end{proof}

\begin{ex}[Moulin.FV.22]{$\star$ Formule des trapèzes}
	Soit $f:[-1,1]\to\C$ une fonction de classe $\mathcal{C}^2$. Montrer que $$\left\lvert\int_{-1}^{1}f(t)\ dt-f(1)-f(-1)\right\rvert\leqslant\frac{2}{3}\underset{x\in[-1,1]}{\sup}\lvert f''(x)\rvert$$
\end{ex}
\begin{proof}
	Faire une première IPP : $$\int_{-1}^{1}f(t)\ dt=\left[tf(t)\right]_{-1}^1-\int_{-1}^{1}tf'(t)\ dt$$ En faire une deuxième avec la bonne fonction pour annuler le crochet, ce qui donne : $$\int_{-1}^{1}f(t)\ dt-f(1)-f(-1)=\underbrace{\left[f'(t)\left(\frac{1}{2}-\frac{t^2}{2}\right)\right]_{-1}^1}_{=0}+\int_{-1}^{1}f''(t)\left(\frac{t^2}{2}-\frac{1}{2}\right)\ dt$$ Ensuite, majorer en valeur absolue et calculer l'intégrale $$\int_{-1}^{1}\left(\frac{1}{2}-\frac{t^2}{2}\right)\ dt=\frac{2}{3}$$
\end{proof}

\begin{ex}[Moulin.FV.23]{Utilisation inhabituelle de Taylor reste intégral}
	Soit $f:\R\to\R$ de classe $\mathcal{C}^{n+1}$. Montrer que $$f(0)=\sum_{k=0}^{n}\frac{(-1)^kx^k}{k!}f^{(k)}(x)+o(x^n)$$ Généraliser ensuite ce résultat au cas où $f$ est seulement supposée de classe $\mathcal{C}^{n}$.
\end{ex}
\begin{proof}
	Dans le cas $\mathcal{C}^{n+1}$, appliquer la formule de Taylor avec reste intégral dans le sens inhabituel. Ensuite, montrer que le reste intégrale est un $o(x^n)$ brutalement. A faire dans le cas $\mathcal{C}^n$.
\end{proof}

\begin{ex}[Moulin.FV.24]{$\star$}
	Soit $f\in\mathcal{C}^5([a,b],\K)$. On pose : $$\Delta=f(b)-f(a)-\frac{b-a}{2}(f'(b)+f'(a))+\frac{(b-a)^2}{12}(f''(b)-f''(a))$$
	\begin{enumerate}
		\item Montrer : $$\lvert\Delta\rvert\leqslant\frac{(b-a)^5}{720}\underset{[a,b]}{\max}\lvert f^{(5)}\rvert$$
		\item Lorsque $\K=\R$, montrer qu'il existe $c\in]a,b[$ tel que : $$\Delta=\frac{(b-a)^5}{720}f^{(5)}(c)$$
	\end{enumerate}
\end{ex}

\begin{ex}[Moulin.FV.25]{}
	Soit $f\in\mathcal{C}^2([a,b],E)$ telle que $f(a)=f(b)=0$. Montrer : $$\forall x\in[a,b],\ \lVert f(x)\rVert\leqslant M\frac{(x-a)(b-x)}{2}\ \ \text{où}\ \ M=\underset{[a,b]}{\sup}\lVert f''\rVert$$
	\begin{enumerate}
		\item dans le cas où $E=\R$ ;
		\item $\star$ dans le cas général.
	\end{enumerate}
\end{ex}

\begin{ex}[Moulin.FV.26]{}
	Soit $f\in\mathcal{C}^2([a,b],\C)$ telle que $f'(a)=f'(b)=0$. Montrer : $$\lvert f(b)-f(a)\rvert\leqslant\frac{(b-a)^2}{4}\underset{[a,b]}{\sup}\lvert f''\rvert$$
\end{ex}
\begin{proof}
	Écrire la formule de Taylor avec reste intégral à l'ordre 2 en $\displaystyle\frac{a+b}{2}$ par rapport à $a$ puis par rapport à $b$. Sommer les deux relations puis majorer et calculer l'intégrale pour obtenir la quantité souhaitée.
\end{proof}

\begin{ex}[Moulin.FV.27]{$\Delta$ Inégalité de Kolmogorov}
	Soit $f$ de classe $\mathcal{C}^2$ sur $\R$. On suppose que $f$ et $f''$ sont bornées respectivement par $M_0$ et $M_2$.
	\begin{enumerate}
		\item Montrer que pour tout $h>0$, la fonction $f'$ est bornée par $\displaystyle\frac{2M_0}{h}+\frac{hM_2}{2}$.
		\item En déduire que $f'$ est bornée par $2\sqrt{M_0M_2}$.
		\item Soit $f$ définie sur $[0,1]$ par : $$\forall x\in[0,1],\ f(x)=2(x-1)^2-1$$ Quelles sont les bornes de $f$, $f'$ et $f''$ sur $[0,1]$ ? En utilisant la fonction cosinus, prolonger $f$ est une fonction de classe $\mathcal{C}^2$ sur $\R_+$ pour laquelle la majoration de la question précédente est optimale.
	\end{enumerate}
\end{ex}

\begin{ex}[Moulin.FV.28]{Égalité de Taylor-Lagrange}
	Soit $f\in\mathcal{C}^n([0,a],\R)$ avec $a>0$.
	\begin{enumerate}
		\item Montrer que pour tout $x\in]0,a[$, il existe $c\in]0,1[$ tel que $$f(x)=\sum_{k=0}^{n-1}\frac{x^k}{k!}f^{(k)}(0)+\frac{x^n}{n!}f^{(n)}(cx)$$ On pourra considérer un réel $A$ tel que $$f(x)=\sum_{k=0}^{n-1}\frac{x^k}{k!}f^{(k)}(0)+\frac{x^n}{n!}A$$
		\item On suppose de plus $f$ de classe $\mathcal{C}^{n+1}$ avec $f^{(n+1)}(0)\ne 0$. Montrer l'unicité d'un tel $c$ pour $x$ assez proche de $0$. On le note $c_x$.
		\item Déterminer la limite de $c_x$ lorsque $x$ tend vers $0$.
	\end{enumerate}
\end{ex}
\begin{proof}
	Avant, la première question était au programme.
	\begin{enumerate}
		\item Suivre l'indication.
		\item A faire...
		\item Peut-être faire un DL dans la première expression ?
	\end{enumerate}
\end{proof}

\begin{ex}[Moulin.FV.29]{$\Delta\star$ Théorème de division}
	Soit $f\in\mathcal{C}^\infty(\R,E)$. Montrer que l'application $g:x\mapsto\displaystyle\frac{f(x)-f(0)}{x}$ peut se prolonger en une fonction de classe $\mathcal{C}^\infty$ sur $\R$.
\end{ex}
\begin{proof}
	Appliquer la formule de Taylor avec reste intégral dans le sens inhabituel et le théorème hors-programme du cours (limite de la dérivée à l'ordre $n$). Sinon, on peut utiliser des intégrales à paramètres, c'est beaucoup plus confortable. En effet : $$g(x)=\int_{0}^{1}f'(tx)\ dt$$ est le seul prolongement convenable. On peut prouver que cette fonction est alors de classe $\mathcal{C}^\infty$ avec les théorèmes d'intégrales à paramètres.
\end{proof}

\begin{ex}[Moulin.FV.30]{$\star$ Racine carrée d'une fonction paire}
	Soit $f:\R\to E$ une fonction paire de classe $\mathcal{C}^\infty$. Montrer qu'il existe $g:\R_+\to E$ de classe $\mathcal{C}^\infty$ telle que $$\forall x\in\R,\ f(x)=g(x^2)$$ On pourra utiliser le théorème de division (exercice précédent).
\end{ex}
\begin{proof}
	A faire...
\end{proof}

\begin{ex}[Moulin.FV.31]{$\Delta$}
	Soit $f\in\mathcal{C}^\infty(\R)$ et $P$ un polynôme de degré impair. On suppose : $$\forall x\in\R,\ \forall n\in\N,\ \lvert f^{(n)}(x)\rvert\leqslant\lvert P(x)\rvert$$ Montrer que $f$ est nulle.
\end{ex}
\begin{proof}
	$p$ possède une racine $a$. Toutes les dérivées de $f$ sont donc nulles en $a$. Appliquer ensuite l'inégalité de Taylor-Lagrange à $x$ par rapport à $a$ en remarquant que toutes les dérivées sont bornées par une même constante par hypothèse.
\end{proof}

\section{Équations fonctionnelles}
\begin{ex}[Moulin.FV.32]{$\Delta$}
	On s'intéresse au fonctions $f:\R\to\C\st$ telle que $$\forall (x,y)\in\R^2,\ f(x+y)=f(x)f(y)$$
	\begin{enumerate}
		\item Déterminer les fonctions dérivables vérifiant cette condition.
		\item $\star$ Déterminer les fonctions continues vérifiant cette condition. On montrera que toute solution est dérivable, en introduisant une primitive $F$ de $f$ et en intégrant par rapport à une variable.
	\end{enumerate}
\end{ex}
\begin{proof}
	La deuxième question est difficile.
	\begin{enumerate}
		\item Dériver par rapport à une variable pour obtenir une équation différentielle linéaire du premier ordre à coefficients constants. On trouve une solution exponentielle.
		\item Intégrer par rapport à l'une des deux variables puis bien choisir remarquer que la primitive ne peut être constante, sans quoi $f$ serait nulle.
	\end{enumerate}
\end{proof}

\begin{ex}[Moulin.FV.33]{$\Delta$Sous-groupes à un paramètre}
	Soit $\varphi:\R\mapsto\mathcal{GL}_n(\C)$ une fonction continue telle que $$\forall (x,y)\in\R^2,\ \varphi(x+y)=\varphi(x)\varphi(y)$$
	\begin{enumerate}
		\item En introduisant une primitive de $\varphi$, montrer que $\varphi$ est dérivable.
		\item Trouver une équation différentielle vérifiée par $\varphi$ et en déduire que $\varphi$ est de classe $\mathcal{C}^\infty$.
	\end{enumerate}
\end{ex}
\begin{proof}
	\begin{enumerate}
		\item Faire comme dans l'exercice précédent en utilisant un DL au voisinage de $0$.
		\item Dériver par rapport à l'une des deux variables. Récurrence pour la classe $\mathcal{C}^\infty$.
	\end{enumerate}
\end{proof}

\begin{ex}[Moulin.FV.34]{}
	On cherche à déterminer les fonctions continues $f:\R\to\R$ telles que $$\forall (x,y)\in\R^2,\ f(x+y)f(x-y)=f^2(x)-f^2(y)$$
	\begin{enumerate}
		\item Montrer qu'une telle fonction est impaire.
		\item $\star\star$ Déterminer toutes les fonctions de classe $\mathcal{C}^2$ vérifiant cette condition.
		\item Montrer que toute fonction continue vérifiant cette condition est de classe $\mathcal{C}^\infty$. On pourra écrire la condition sous la forme $f(u)f(v)=\ldots$ et primitiver par rapport à $v$.
	\end{enumerate}
\end{ex}
\begin{proof}
	La deuxième question est très difficile.
	\begin{enumerate}
		\item $x=y=0$ donne $f(0)=0$. Ensuite, jouer avec $x=0$ ou $x=y$ ou $x=-y$.
		\item A faire...
		\item Écrire la relation sous la forme suggérée et dériver.
	\end{enumerate}
\end{proof}

\begin{ex}[Moulin.FV.35]{$\star$}
	Déterminer les fonctions $f\in\mathcal{C}^0(\R,\R)$ telle que $$\forall (x,y)\in\R^2,\ f(x+y)+f(x-y)=2f(x)f(y)$$ On pourra s'inspirer des exercices précédents.
\end{ex}
\begin{proof}
	Cet exercice est difficile. On montre avec $x=y=0$ que $f(0)$ vaut $0$ ou $1$. Si $f(0)=0$, on a $f=0$ en prenant $y=0$. On suppose donc $f(0)=1$. Ensuite, on primitive pour montrer que $f$ est dérivable. On se ramène à une équation d'ordre $2$. On trouve des solutions en $\cos$ et $\cosh$.
\end{proof}

\end{document}